\section{Заключение}

В работе сформулирована вероятностная постановка задачи выбора времени остановки при рекурсивном прогнозировании речи. Речь рассматривается как случайная последовательность на конечном словаре, а рекурсивная модель строит несколько последовательных распределений следующей языковой единицы. По полной последовательности потерь определяется оракульное время остановки, минимизирующее сумму ошибки и стоимости вычислений.

Основная статистическая задача состоит в оценке условного распределения этого оракульного времени по информации, доступной на текущем внутреннем шаге. Правило остановки отделено от модели распределения. Это позволяет учитывать не только вероятность остановки сейчас, но и вероятность того, что оптимальный момент уже был в прошлом или еще находится впереди.

Рассмотрены шесть семейств распределений на \(\{0,\ldots,K\}\): конечное дискретное, сглаженная целевая разметка по потерям, смесь геометрических, дискретизированная смесь экспоненциальных, степенной спад на отрезке и отрицательное биномиальное распределение. Для принятия решения используются пять правил: накопленная вероятность, малая вероятность будущего улучшения, условная интенсивность по предсказанной мере, квантиль и бюджетное правило.

На данных ARC-AGI приведена сводка валидационных экспериментов по сетке стратегий остановки, эмпирическое распределение оракульного \(\tau^\star\) и сравнение базового обучения с дообучением только головы оракула. Для языкового трека зафиксирована иллюстративная постановка на публичных корпусах с плейсхолдерными численными значениями.

Направления дальнейшей работы: (i) полный языковой свип с тем же протоколом оценки; (ii) формальные оценки дрифта \(\widehat\nu_{\mathrm{train}}\) vs \(\widehat\nu_{\mathrm{val}}\) для \(\tau^\star\); (iii) построение калибровочных кривых по полному массиву \texttt{stopping\_eval.json}.
